\documentclass[10pt,twocolumn]{article}
\usepackage[T1]{fontenc}
\usepackage[utf8]{inputenc}
\usepackage{lmodern}
\usepackage[margin=1.8cm,top=2.4cm,bottom=2cm,columnsep=0.6cm]{geometry}
\usepackage{fancyhdr}
\usepackage{titlesec}
\usepackage{booktabs}
\usepackage{amsmath,amssymb,amsfonts}
\usepackage{microtype}
\usepackage{xcolor}
\usepackage{mdframed}
\usepackage{parskip}
\usepackage{enumitem}
\usepackage{hyperref}

\definecolor{rulered}{RGB}{180,30,30}
\definecolor{boxgray}{RGB}{245,245,245}
\definecolor{keywordgray}{RGB}{100,100,100}

\pagestyle{fancy}
\fancyhf{}
\fancyhead[L]{\textsc{\small Claude Mag}}
\fancyhead[C]{\small\textit{Machine Learning Research Digest}}
\fancyhead[R]{\small 2026-08-17}
\fancyfoot[C]{\small\thepage}
\renewcommand{\headrulewidth}{0.8pt}

\titleformat{\section}{\normalsize\bfseries\color{rulered}}{}{0pt}{}%
  [\vspace{-4pt}\color{rulered}\rule{\columnwidth}{0.4pt}]
\titlespacing{\section}{0pt}{8pt}{4pt}

\hypersetup{colorlinks=true,urlcolor=rulered,linkcolor=black}

\begin{document}
\setlength{\parindent}{0pt}
\setlength{\parskip}{4pt}

{\small\color{keywordgray}\textbf{Keywords:}
  RLHF fairness \textbullet{}
  preference heterogeneity \textbullet{}
  reward modeling \textbullet{}
  alignment}\par\vspace{4pt}

{\LARGE\bfseries\raggedright
  Procedural Fairness Failures in RLHF
  from Preference Averaging}\par\vspace{4pt}

{\large\itshape\raggedright
  Standard RLHF Violates Procedural Fairness by Drowning Minority Preference Signals
  in Aggregated Reward Models; PA-RLHF Separates Optimization Across Preference Modes}\par\vspace{6pt}

{\small\textbf{By} M P V S Gopinadh, Karthik Kamuju, Kummari Avinash,
  John Joshua, Srinivasa Raju Rudraraju}\par\vspace{2pt}

{\small\color{keywordgray}arXiv:2608.10126 \textbullet{}
  ICLR 2026 Workshop on Algorithmic Fairness Across Alignment Procedures (AFAA)}
\par\vspace{2pt}\noindent{\color{rulered}\rule{\columnwidth}{0.6pt}}\vspace{6pt}

Reinforcement Learning from Human Feedback (RLHF) aggregates heterogeneous annotator
preferences into a single reward model, treating all annotators as draws from one
distribution. A new paper shows this induces a \emph{procedural fairness failure}:
majority preference groups dominate reward learning while minority preferences are
systematically under-represented. Preference-Aware RLHF (PA-RLHF) remedies this by
separating reward optimization across preference modes at the reward learning stage.

\begin{mdframed}[backgroundcolor=boxgray,linecolor=rulered,linewidth=1pt,
    innerleftmargin=6pt,innerrightmargin=6pt,
    innertopmargin=6pt,innerbottommargin=6pt]
\textbf{\small AT A GLANCE}\par\vspace{4pt}
\begin{description}[leftmargin=0pt,labelsep=4pt,itemsep=2pt,parsep=0pt,
    font=\small\bfseries]
  \item[Problem:] Standard RLHF averages heterogeneous preferences into a single reward
    model, causing minority preference groups to be under-represented.
  \item[Why it matters:] As LLMs are deployed across diverse populations, preference
    homogeneity is a poor assumption — and its failure is invisible to standard metrics.
  \item[Method:] Define procedural fairness for alignment; prove standard RLHF violates it;
    propose PA-RLHF with per-group reward models.
  \item[Result:] PA-RLHF improves procedural fairness metrics while maintaining outcome
    quality (win rates, helpfulness) across all tested scenarios.
  \item[Evidence:] Gains largest when inter-group preference divergence and population
    size imbalance are greatest.
\end{description}
\end{mdframed}

\section{The Big Picture}

Human preferences are not homogeneous. Annotators employed in RLHF pipelines come from
different cultural backgrounds, hold different values, and have different conceptions
of what constitutes helpful, harmless, and honest AI behavior. Standard RLHF treats all
annotator preferences as if they were noisy draws from a single underlying distribution,
fitting a single aggregated reward model $r_\theta$ to the pooled data.

When preferences are genuinely multi-modal — when subgroups of annotators hold coherently
different values — this aggregation is not merely imprecise. It actively \emph{dilutes}
minority preference signals in proportion to the size of the minority. Smaller groups
contribute less to the aggregate loss, so the reward model learns to satisfy the majority
and discard the minority. The paper calls this a \textbf{procedural fairness failure}:
the process is unfair to minority groups, independent of whether the resulting policy
happens to satisfy them by coincidence.

\section{Why Existing Approaches Fall Short}

Prior work on fairness in RLHF has focused on \textbf{outcome fairness}: ensuring the
final policy produces acceptable outputs for all subgroups. But outcome fairness can
be achieved by coincidence (a policy that satisfies the majority may also happen to
satisfy the minority on some tasks) and can be gamed by optimizing for population-level
statistics. Outcome fairness does not address whether the reward learning \emph{process}
treats all groups equitably.

Existing personalized preference methods (per-user reward models, preference
embeddings) require per-user data at inference time and scale poorly. PA-RLHF
operates at the group level during reward training, requiring no per-user data
after deployment.

\section{Core Mechanism}

Let $\mathcal{A}$ be the annotator pool, partitioned into $K$ preference groups
$G_1, \ldots, G_K$ with group-specific preference distributions
$\mathcal{D}_1, \ldots, \mathcal{D}_K$. The key steps of PA-RLHF are:

\begin{enumerate}[itemsep=2pt,parsep=0pt]
  \item \textbf{Preference clustering:} Identify latent groups $G_k$ using embedding-based
    clustering of annotator preference vectors or label co-occurrence patterns.
  \item \textbf{Per-group reward modeling:} Fit a separate reward model $r_{\theta_k}$
    for each group $G_k$ using only that group's preference data.
  \item \textbf{Fairness-weighted policy optimization:} Combine per-group reward signals
    during PPO training using a fairness weighting $w_k$ that does not simply
    equal the group's population share $\alpha_k = |G_k|/|\mathcal{A}|$.
\end{enumerate}

The fairness weighting $w_k$ is the key design choice: uniform weighting ($w_k = 1/K$)
maximizes procedural equity; intermediate weightings balance procedural and outcome
fairness objectives.

\section{Technical Formulation}

\textbf{Standard RLHF} fits:
\[
r_\theta = \arg\min_\theta \mathbb{E}_{(x,y_w,y_l) \sim \mathcal{D}}\!\left[-\log\,\sigma(r_\theta(x,y_w) - r_\theta(x,y_l))\right]
\]
where $\mathcal{D} = \sum_k \alpha_k \mathcal{D}_k$ is the mixture distribution.

\textbf{PA-RLHF} fits $K$ per-group reward models:
\[
r_{\theta_k} = \arg\min_{\theta_k} \mathbb{E}_{(x,y_w,y_l) \sim \mathcal{D}_k}\!\left[-\log\,\sigma(r_{\theta_k}(x,y_w) - r_{\theta_k}(x,y_l))\right]
\]
and optimizes the policy with the fairness-weighted reward:
\[
r_{\mathrm{fair}}(x,y) = \sum_{k=1}^K w_k \cdot r_{\theta_k}(x,y), \quad \sum_k w_k = 1
\]

\textbf{Procedural fairness} is defined as: the gradient contribution of group $G_k$
to the aggregate reward model should be proportional to its ``representation right''
rather than its population size. Formally, standard RLHF satisfies this only when
$\alpha_k = 1/K$ (equal group sizes), which is rarely the case in practice.

\section{Learning or Inference Procedure}

PA-RLHF requires one additional step versus standard RLHF:
\begin{enumerate}[itemsep=2pt,parsep=0pt]
  \item \textbf{Cluster annotators} into $K$ groups (one-time offline step).
  \item \textbf{Train $K$ reward models} in parallel (same cost as standard training
    when using parameter-efficient methods; $K \times$ cost otherwise).
  \item \textbf{Run PPO} with $r_{\mathrm{fair}}$ as the reward signal.
\end{enumerate}
Inference is identical to standard RLHF — no per-user data is required.

\section{What the Guarantee Says}

The paper proves that standard RLHF is procedurally unfair whenever annotator
group sizes differ ($\alpha_k \neq 1/K$) and group preferences are non-identical.
PA-RLHF with uniform weighting ($w_k = 1/K$) is the unique solution that achieves
procedural fairness under this definition. No guarantee is provided for outcome
fairness — PA-RLHF may improve or maintain outcome metrics, but cannot guarantee it.

\section{Experimental Findings}

\begin{itemize}[itemsep=2pt,parsep=0pt]
  \item \textbf{Procedural fairness:} PA-RLHF improves procedural fairness metrics
    across all tested scenarios with heterogeneous annotator groups.
  \item \textbf{Outcome quality:} Win rates and helpfulness metrics are maintained
    or slightly improved relative to the standard RLHF baseline.
  \item \textbf{Group imbalance sensitivity:} Gains grow with group size imbalance
    ratio — the regime most common in real crowdsourced annotation datasets.
  \item \textbf{Preference divergence sensitivity:} Gains also grow with inter-group
    preference divergence, confirming the mechanism targets the right failure mode.
\end{itemize}

\section{Ablations and Interpretation}

\textbf{How many groups $K$?} The paper tests $K \in \{2, 3, 5\}$; fairness gains
are present for all values, with $K=3$ providing the best balance of fairness and
computation overhead.

\textbf{Fairness weighting $w_k$:} Uniform ($w_k = 1/K$) maximizes procedural
fairness but can hurt outcome quality in extreme scenarios. Population-proportional
($w_k = \alpha_k$) recovers standard RLHF. An intermediate value provides the
best empirical outcome-fairness tradeoff.

\textbf{What about outcome fairness?} The paper notes that procedural fairness and
outcome fairness are not equivalent: PA-RLHF improves procedural fairness without
guaranteeing outcome fairness, and a system can achieve outcome fairness while
violating procedural fairness. Both matter and both need to be measured.

\section*{Reference}

M P V S Gopinadh, Karthik Kamuju, Kummari Avinash, John Joshua,
Srinivasa Raju Rudraraju.
\textbf{Procedural Fairness Failures in RLHF from Preference Averaging}.
arXiv:2608.10126, August 2026.
ICLR 2026 Workshop on Algorithmic Fairness Across Alignment Procedures and Agentic Systems.
\url{https://arxiv.org/abs/2608.10126}

\end{document}
